\subsubsection{An exact scalar factory}

From here on fix the absolute constants
\[
 \eta=1/64,\qquad \alpha=65/64,\qquad
 \varphi(z)=\sqrt{\alpha-z^2},\qquad
 c_0=17/16,\qquad M_0=2^{22},\qquad J_0=2^{25}.
\]
Consider a row with \(0<\rho\le1/32\) and \(0\le a\le4\rho\).
Negative \(a\) is handled by complementing source and output signs.
Write
\[
 q=\tanh\rho,\quad t=\tanh a,\quad c=tq,\quad
 G(z)=\tanh(\rho z)/q,\quad
 F(z)=\frac{G(z)+c}{1+cG(z)}.
\]
The function
\[
 H(z)=\frac{F(z)-z}{1-z^2}=\sum_{k\ge0}h_kz^k
\]
has removable singularities at \(z=\pm1\).

\begin{lemma}[Coefficients of the correction]\label{lem:fusedcoeff}
The function \(H\) is analytic on a neighborhood of the closed
complex disk \(|z|\le2\). Its first coefficients obey
\[
 h_0=c\ge0,\qquad h_1=(1-c^2)\rho/q-1\ge0,
 \qquad h_0\le4\rho^2,\quad h_1\le\rho^2/3.
\]
For \(k\ge2\),
\begin{equation}\label{eq:fusedtail}
 |h_k|\le200\rho^4\,2^{-k}.
\end{equation}
Also, with \(u=t/\rho\in[0,4]\),
\[
 |h_0-u\rho^2|\le2\rho^4,
 \qquad |h_1-\rho^2/3|\le100\rho^4.
\]
\end{lemma}
\paragraph{Proof idea.}
Expand the scalar correction on a complex disk where all denominators stay
away from zero. Cauchy's estimate bounds the entire coefficient tail, while
direct differentiation identifies the two coefficients handled separately by
the factory.

\begin{proof}
We use the complex bound \eqref{eq:tanhfifth}.
On \(|z|=2\), use \(q/\rho\ge1-\rho^2/3\) to obtain
\[
 \left|G(z)-z-\frac{\rho^2}3z(1-z^2)\right|
 \le200\rho^4,\quad |G(z)-z|\le4\rho^2,\quad |G(z)|<3.
\]
Here the two fifth-order Taylor remainders, after division by
\(\rho\), total at most \(170\rho^4\); the denominator and
the quadratic correction increase this to less than
\(200\rho^4\).  The denominator \(1+cG\) has modulus at
least \(1-12\rho^2>49/50\). The sharper numerator bound
\(|G+c|\le2+8\rho^2\) also gives
\[
 |F(z)|\le\frac{2+8\rho^2}{1-12\rho^2}<3
 \qquad(|z|\le2).
\]
Moreover
\[
 F-z-\frac{\rho^2}3z(1-z^2)-c(1-z^2)
 =G-z-\frac{\rho^2}3z(1-z^2)
 +\frac{c\{z^2-G^2-cG(1-z^2)\}}{1+cG}.
\]
The last fraction has modulus at most
\(4\rho^2(20\rho^2+60\rho^2)/(49/50)\).
The entire expression is therefore bounded by \(600\rho^4\).
On this circle \(|1-z^2|\ge3\), so
\begin{equation}\label{eq:fusedHbound}
 |H(z)-c-(\rho^2/3)z|\le200\rho^4.
\end{equation}
The denominators above are nonzero throughout the disk.  Since
\(F(1)=1\) and \(F(-1)=-1\), its apparent singularities in
\(H\) are removable.  Cauchy's estimate now proves
\eqref{eq:fusedtail} and the assertion about \(h_1\).

Direct differentiation at zero gives the displayed values of
\(h_0,h_1\).  The real Taylor bound implies
\(q\le\rho(1-\rho^2/4)\), whereas
\(\rho/q\le1+\rho^2/3\) by \eqref{eq:xcothx}. Consequently
\[
 h_1\ge\rho^2/4-17\rho^4>0,
\]
using \(c^2\le16\rho^4\) and \(\rho/q\le17/16\).
Finally \(|tq-t\rho|\le4\rho^4/3\), proving the remaining
estimate.
\end{proof}

Here is an explicit factory for \(F\).  Read one source sign
\(X_1\).  Put \(A=h_0\), \(B=h_1\), and \(e_0=200\rho^4\).
After \(X_1=+1\), choose a majority-of-three continuation with
probability \(2(B-A)_+\).  After \(X_1=-1\), choose such a
continuation with probability \(2B+2\min(A,B)\), or an OR
continuation with probability \(2(A-B)_+\).  A majority
continuation stops as soon as two signs agree.  An OR continuation
after the negative first sign reads one further sign and returns it.

In either case, reserve the additional probability \(e_0\) for a
tail proposal.  Otherwise return \(X_1\).  For a tail proposal,
draw \(K\ge2\) with \(\Pr(K=k)=2^{1-k}\), and choose its
positive or negative correction with respective probabilities
\[
 \frac{(h_k)_+}{200\rho^4 2^{-k}},\qquad
 \frac{(-h_k)_+}{200\rho^4 2^{-k}}.
\]
On rejection return \(X_1\).  On acceptance, read \(X_2\).
If \(X_1=X_2\), return their common sign.  Otherwise return
the chosen sign times the product of \(k\) fresh source signs.
All category probabilities are valid: their total is at most
\(9\rho^2+200\rho^4<1/100\).  The positive-part formulas
avoid deciding an exact equality between computable real numbers.

\begin{lemma}[Exactness and source cost]\label{lem:fusedfactory}
The factory terminates almost surely and has mean \(F(z)\).
Define
\[
 E(A,B,z)=B(3-z^2)-2Az+(A-B)_+(1+z).
\]
Its expected number of source signs is
\[
 A_{\mathrm f}(\rho,t,z)
 =1+E(h_0,h_1,z)
   +\sum_{k\ge2}|h_k|\{2+k(1-z^2)\}.
\]
In particular,
\begin{equation}\label{eq:fusedarity}
 A_{\mathrm f}\le1+\rho^2 e(u,z)+1000\rho^4,
\end{equation}
where
\[
 e(u,z)=
 \begin{cases}
  1-z^2/3-2uz,&0\le u\le1/3,\\
  2/3-z/3-z^2/3+u(1-z),&u\ge1/3.
 \end{cases}
\]
\end{lemma}
\paragraph{Proof idea.}
Condition on the first source sign to compute the majority and OR
corrections. Opposite first signs expose the product correction, whose
geometric mixture supplies the remaining Taylor series. The same
conditioning counts every additional source request.

\begin{proof}
Conditioning on \(X_1\), the majority and OR continuations
change its mean by precisely \((1-z^2)(A+Bz)\).
The signed tail component has mean
\(z\pm(1-z^2)z^k/2\), since two source signs disagree with
probability \((1-z^2)/2\).  Its unconditional probability is
\(2(h_k)_+\) or \(2(-h_k)_+\), independently of \(X_1\).
Absolute convergence therefore gives mean
\(z+(1-z^2)H(z)=F(z)\).

Conditional on a positive first sign, a majority continuation
uses \((3-z)/2\) additional signs on average; after a negative
one it uses \((3+z)/2\).  The OR continuation uses one.
This yields \(E(A,B,z)\).  A tail component uses, beyond the
first sign, \(1+k(1-z^2)/2\) signs on average, proving the
exact cost formula.  Its geometric index is finite almost surely.

The function \(E\) is continuous, with separate Lipschitz
constants \(2\) in \(A\) and \(3\) in \(B\).  Hence the first
two coefficient errors cost at most \(304\rho^4\).
The tail costs at most
\[
 200\rho^4\sum_{k\ge2}(k+2)2^{-k}=500\rho^4.
\]
Substitution of \((A,B)=(u\rho^2,\rho^2/3)\) gives
\eqref{eq:fusedarity}.  Known scalar decisions terminate almost
surely by the same interval-comparison lemma as the original
factory.
\end{proof}

\subsubsection{The weighted estimate}

Let \(k=(1-t^2)/(1-t^2q^2)\).  After the internal gate, the
normalized mean before the center error is \(v=kF(z)\).
As in Section~\ref{app:potential}, concavity gives
\(\sum_jp_j\varphi(v_j)\le\varphi(z)\), and the ratio
\(\theta/\varphi(\theta v)\) increases with \(\theta\).
It therefore suffices to bound
\[
 W_{\mathrm f}=A_{\mathrm f}k\frac{\varphi(z)}{\varphi(kF(z))}.
\]

\begin{lemma}[A fixed potential bound]\label{lem:fusedlocal}
For \(0<\rho\le1/32\) and \(0\le a\le4\rho\),
\begin{equation}\label{eq:fusedlocal}
 \log W_{\mathrm f}\le\frac{17}{16}\rho^2+2^{22}\rho^4.
\end{equation}
For \(a>4\rho\), the original majority and race factory obeys
the same weighted estimate.  Negative biases are complemented.
The center error multiplies the estimate by
\(\exp(64\varepsilon/r_\ell)\).
\end{lemma}
\paragraph{Proof idea.}
Combine the local source expansion with the logarithm of the potential
ratio. Its leading quadratic form can be maximized explicitly over the
normalized bias; a uniform fourth-order remainder controls the error.

\begin{proof}
Put \(d=\alpha-z^2\), and
\[
 b(u,z)=(1-z^2)(z/3+u)-u^2z.
\]
The bounds in Lemma~\ref{lem:fusedcoeff} imply
\[
 |kF(z)-z-\rho^2b(u,z)|\le1000\rho^4,
 \qquad |kF(z)-z|\le22\rho^2.
\]
For completeness, the first remainder is at most
\(602\rho^4+80\rho^4+17\rho^4<700\rho^4\):
these are respectively the remainder in \(F-z\), the term
\(t^2(F-z)\), and the difference \(k-(1-t^2)\).
Here \(|F-z|\le5\rho^2\), \(|b|<21\), and
\(0\le k-(1-t^2)\le17\rho^4\).
Also \(\log k\le-u^2\rho^2+17\rho^4\).
Taylor's theorem for \(\log\varphi\), whose first and second
derivatives have absolute bounds \(1/\eta\) and
\(9/(4\eta^2)\), gives
\[
 \log W_{\mathrm f}
 \le C_\eta(u,z)\rho^2+
 \left(1017+\frac{1000}\eta+
                   \frac{1089}{2\eta^2}\right)\rho^4,
 \qquad
 C_\eta=e(u,z)-u^2+\frac{zb(u,z)}d.
\]
At \(\eta=1/64\), the coefficient of \(\rho^4\) is
\(2295289<2^{22}\).

We show \(C_\eta\le17/16\) throughout \(0<\eta\le1/64\).
For \(u\le1/3\),
\[
 C_\eta=1-\frac{\eta z^2}{3d}
       -z\left(1+\frac\eta d\right)u-\frac\alpha d u^2.
\]
If \(z\ge0\), this is at most one.  Otherwise maximize the
quadratic over all real \(u\).  With \(y=z^2\), the desired
bound is equivalent to
\[
 P_\eta(y)=3(1-y)(2y-1)^2
       +\eta(6-35y+48y^2)+\eta^2(3-32y)\ge0.
\]
For \(y=1/2+v\), \(0\le v\le1/2\), the first term is
nonnegative, the coefficient of \(\eta\) is
\(1/2+13v+48v^2\ge1/2\), and that of \(\eta^2\) is at
least \(-29\).  Thus \(P_\eta\ge\eta/2-29\eta^2>0\).
For \(y=1/2-v\), \(0\le v\le1/2\),
\[
 P_\eta\ge6v^2-13\eta v+\eta/2-13\eta^2
 =6(v-13\eta/12)^2+\eta/2-481\eta^2/24>0.
\]
Both inequalities use only \(\eta\le1/64\).

For \(u\ge1/3\),
\[
 C_\eta=2/3-z/3-\frac{\eta z^2}{3d}
       +\left(1-\frac{\eta z}d\right)u-\frac\alpha d u^2.
\]
If its unconstrained maximum occurs below \(1/3\), the maximum
on the half-line is attained at the boundary, already covered.
Otherwise \(z^2+\eta z\le\alpha/3\), implying \(z\ge-3/5\).
Completing the square and discarding the nonpositive term
\(\eta z^2(\eta/(4\alpha)-1/3)/d\) gives
\[
 C_\eta\le11/12-
      \left(1/3+\frac\eta{2\alpha}\right)z-
      \frac{z^2}{4\alpha}.
\]
This quadratic decreases for \(z\ge-3/5\).  Its value at
\(-3/5\) is
\(67/60+(30\eta-9)/(100\alpha)\le21/20<17/16\), since
\(\eta\le1/64<7/110\).  This proves the small-bias assertion.

For \(a>4\rho\), monotonicity and
\(\tanh(4\rho)\ge4\rho-(4\rho)^3/3\) give
\(t\ge(7/2)q\).  Lemma~\ref{lem:potential-majority-bias} gives
\[
 B\le\frac{-(1-2q^2)t^2+3tq}{1-q^2}\le0,
\]
because \(-1/2+7q^2<0\).  Its remaining majority contribution
is at most \(\rho^2+4\rho^4/\eta^2\), covered by
\eqref{eq:fusedlocal}.  The center shift and the factor
\(\theta=q/r_\ell\le1\) are treated exactly as in
\eqref{eq:potentialW}.
\end{proof}

\subsubsection{Consequences for depth}

Use the original factory at levels \(\ell\le J_0\), and the
new scalar factory above \(J_0\) when \(\rho\le1/32\) and
\(|a|\le4\rho\).  All tests use stored rational quantities.
Outside this case retain the original factory.

\begin{theorem}[Improved critical bound]\label{thm:fusedcritical}%
\label{thm:criticalrefined}
For row norms at most one, this sampler is exact and satisfies
\[
 \E N_{\rm calls}=O((D+1)^{35/32}).
\]
Its preprocessing obeys the original bound, and its expected
online bit work is \(O(\bits^4(D+1)^{35/32})\).
Preprocessing may select between this bound and the quadratic
critical bound of Theorem~\ref{thm:upper}.
\end{theorem}
\paragraph{Proof idea.}
Use the fused estimate only above a fixed cutoff. Weighted offspring
products across those layers give a power of depth, and the final output
gate reduces that power by one half. The original sampler bounds the finite
prefix.

\begin{proof}
When \(D>J_0\), every squared row radius above level \(J_0\)
is at most \(2/J_0<1/1024\).  On this interval the derivative
of \(c_0v+M_0v^2\) is at most
\(17/16+4M_0/J_0=25/16<2\).
The cumulative rounding and center terms are each less than
\(1/100\), as in Section~\ref{app:potential}.
The ideal fourth-power contribution to a suffix starting at
\(k\ge J_0\) is at most
\[
 \frac{9M_0}{4(k+1)}\le\frac9{32}<1.
\]
Consequently its weighted offspring product is at most
\[
 e^2\left(\frac{D+1}{k+1}\right)^{51/32}.
\]
The original normalized subtree bound
\(36(J_0+1)^{5/2}\) applies below the cutoff because those
levels retain the original factory.  Keep the root factor
\(r_D\le2/\sqrt{D+1}\).  Since
\(\max\varphi/\min\varphi=\sqrt{65}<9\),
\(e^2<8\), and \(\sum_{j\ge1}j^{-51/32}<3\),
\[
 \E N_{\rm calls}
 \le5617(J_0+1)^{29/32}(D+1)^{35/32}
 <10^{11}(D+1)^{35/32}.
\]
The final inequality uses \(J_0+1\le2^{26}\), so
\((J_0+1)^{29/32}<2^{24}\).
For \(D\le J_0\), the original \(80(D+1)^2\) bound is
smaller than the same displayed estimate.
Conditional charging and almost-sure termination are unchanged.
\end{proof}

\begin{proposition}[Improved fixed supercritical rate]
\label{prop:fusedrate}
For \(s=1+\delta\), \(0<\delta\le2^{-16}\), let
\(r_*\) be the positive fixed point of \(r=\tanh(sr)\), and
put $u_*=sr_*$. There is an explicitly computable rate
$\Gamma_s=(17/16)u_*^2+O(u_*^4)$ such that, for every network
with actual maximum row norm $S\le s$, the sampler with reference
gain $\max\{1,S\}$ obeys
\[
 \E N_{\rm calls}=O_s(e^{D\Gamma_s}),\qquad
 \frac{51}{16}\delta\le\Gamma_s
 \le\frac{51}{16}\delta+O(\delta^2).
\]
The constants in the rate estimates are absolute; explicit choices
for the rate and prefactor are given in the proof.
The prefactor is for fixed \(s\); the original uniform
transition-window theorem remains in force.
\end{proposition}
\paragraph{Proof idea.}
Compare all sufficiently late radii with the positive fixed point.
Contraction bounds the accumulated transient error, so the limiting weighted
offspring factor determines the exponential rate.

\begin{proof}
Use
\[
 \Gamma_s=\frac{17}{16}u_*^2+2^{22}u_*^4,\qquad H_0=2^{14},
\]
and the computable prefactor
\[
 C_\delta=1+\sqrt{65}
 \exp\left(1+\frac{3H_0}{200}+\frac{8H_0}{\delta}\right)
 \left(\exp(23(J_0+1))+
                  \frac1{1-\exp(-51\delta/16)}\right).
\]
For \(D>J_0\) and \(j\ge J_0\), the comparison radii satisfy
\[
 s^2r_j^2\le4\delta+6/J_0+24/(100J_0)<1/1024.
\]
On \([0,1/1024]\), the derivative of
\(17v/16+2^{22}v^2\) is at most
\(17/16+8192<H_0\).
Write \(\lambda=s(1-r_*^2)\).  The ideal-radius contraction
and the rounding proof give the suffix bound
\[
 \exp\left(1+3H_0/200+8H_0/\delta\right)
 \exp((D-k)\Gamma_s).
\]
Here \(1-\lambda\ge\delta\), as in
Lemma~\ref{lem:nearcriticalgap}.  Below \(J_0\), the
original normalized subtree has mean at most
\(\exp(23(J_0+1))\).  Weighted charging and a geometric sum
prove the stated bound, including \(D\le J_0\).
Finally
\(3\delta\le u_*^2\le3\delta+(3/2)\delta^2\), and
\(3+(3/2)\delta\le25/8\).  Thus the quadratic error is at
most
\[
 \left(51/32+2^{22}(25/8)^2\right)\delta^2
 <40960002\delta^2.
\]
\end{proof}

\subsubsection{Finite-bit implementation}

Only the small-radius branch is new.  The probabilities
\(h_0,h_1\), their positive parts, and their minimum are
computable with the existing scalar interval oracles; all these
operations are Lipschitz.  No sign or equality oracle is required.
For the tail coefficients there is a direct formal-series oracle.
Set \(A=\rho/q\) and
\[
 X(z)=\frac{e^{2\rho z}-1}{2\rho}
      =\sum_{k\ge1}\frac{(2\rho)^{k-1}}{k!}z^k.
\]
Then
\[
 F(z)=\frac{c+(A+c\rho)X(z)}{1+(\rho+cA)X(z)}.
\]
Its denominator has constant coefficient one.  Truncating at
degree \(K\), applying the quotient recurrence, and then
dividing \(F-z\) formally by \(1-z^2\) computes \(h_K\).
The computation rounds after each arithmetic operation, as in
Lemma~\ref{lem:coefficient-four}; it never carries a growing
common rational denominator.

Here are explicit error bounds for this new quotient.  Let
\(\zeta=2^{-V}\) be the rounding unit and first enclose \(A,c\)
to absolute error \(\zeta\).  Computing \(A\) needs only
\(O(\bits)\) additional guard bits because \(q\ge\rho/2\)
and a nonzero dyadic row has \(\log(1/\rho)=O(\bits)\).
The recurrence
\(x_1=1\), \(x_{j+1}=2\rho x_j/(j+1)\)
has multipliers at most \(1/32\), so each computed \(x_j\)
has error at most \(2\zeta\).  The numerator and denominator
coefficients of the displayed quotient then have error at most
\(10\zeta\) apiece.  Its true nonconstant denominator norm
is less than \(1/8\):
\[
 \rho+cA<1/16,\qquad \sum_{j\ge1}|x_j|<2.
\]
Also \(|F(z)|<3\) on \(|z|=2\), so every true coefficient
of \(F\) has absolute value at most three.

If the previous coefficient errors are at most \(E\le1\),
one quotient step has error at most
\(E/8+(50K+20)\zeta\), including the rounded products and
sum.  Induction bounds every error through degree \(K\) by
\(100(K+1)\zeta\).  The final recurrence
\(h_j=F_j-\mathbf1_{\{j=1\}}+h_{j-2}\), with missing
negative-index terms zero, gives error at most
\(101(K+1)^2\zeta\).  Taking
\(V=m+4\bits+4K+64\), and replacing \(m\) by \(m+2\)
for an enclosure of width \(2^{-m}\), validates the preceding
induction and supplies the requested certificate.
There are \(O(K^2)\) operations on
\(O(\bits+K+m)\)-bit integers.  Schoolbook multiplication and
division therefore give \(O((\bits+K+m)^4)\) bit work.
The acceptance denominator \(200\rho^4 2^{-K}\) requires only
\(O(\bits+K)\) further precision bits, preserving that bound.

The geometric tail proposal gives finite moments of every order
for \(k\).  Interval comparisons have geometric precision tails.
The conditional distribution of a coefficient sign is implemented
through its two nonnegative parts in the displayed probabilities.
An accepted branch never requires resolving whether a zero
coefficient is positive.  The original reached-trial lemma charges
all arithmetic, rejected proposals, and subsequent source work.
Thus a reached invocation costs \(O(\bits^4)\) expected scalar
bit work, including every precision refinement.  No table of
coefficients is added to preprocessing storage.
